{
\section*{Resumen}\label{resumen}
\addcontentsline{toc}{chapter}{Resumen}
\thispagestyle{empty}
\let\clearpage\relax

En el presente documento se realizó el diseño de un biorreactor tipo airlift para el
cultivo de hongos coprófilos pertenecientes al filo Ascomycota, con el objetivo de
proporcionar una plataforma experimental confiable y reproducible. El sistema
contará con monitoreo y control de los parámetros: disponibilidad de nutrientes,
temperatura, presión y nivel, permitiendo al investigador caracterizar el
comportamiento del organismo mientras que asegura la estabilidad de las
condiciones de experimentación. De igual manera, se integrará un sistema de
monitoreo de las variables iluminación y oxigenación más no se controlarán. A lo
largo de este documento se presenta el proceso de diseño del contenedor de los
ascomicetos, consideraciones para la selección de materiales y componentes, así
como simulaciones y cálculos que validan la viabilidad de su construcción. A su
vez, se encuentran contenidos planos, diagramas eléctricos, bosquejos de la
interfaz, sistemas de control y plan de manufactura.

\subsection*{Palabras clave}\label{palabras}
\addcontentsline{toc}{chapter}{Palabras clave}

Biorreactor airlift, ascomicetos coprófilos, sistema controlado, monitoreo automatizado.

\newpage
\section*{Abstract}\label{abstract}
\addcontentsline{toc}{chapter}{Abstract}

This document presents the design of an airlift bioreactor for the cultivation of
coprophilous fungi belonging to the Ascomycota phylum, aiming to provide a reliable
and reproducible experimental platform. The system features monitoring and control
of key parameters: nutrient availability, temperature, pressure, and level, allowing
the researcher to characterize the organism's behavior while ensuring stable
experimental conditions. Additionally, a monitoring system for lighting and
oxygenation variables is integrated, though these will not be controlled. Throughout
this document, the design process for the ascomycete containment tank is detailed,
including material and component selection criteria, as well as simulations and
calculations that validate the feasibility of its construction. Furthermore, it contains
blueprints, electrical diagrams, interface sketches, control systems, and a
manufacturing plan.

\subsection*{Keywords}\label{keywords}
\addcontentsline{toc}{chapter}{Keywords}

Airlift bioreactor, coprophilous ascomycetes, controlled system, automated monitoring.
}
